\documentclass[10pt,a4paper]{amsart}
\usepackage[margin=19mm]{geometry}
\usepackage{amsmath,amssymb,mathtools}
\usepackage[scr=rsfs,cal=euler]{mathalfa}
\usepackage{xfrac}
\usepackage[unicode,hidelinks,bookmarks=false]{hyperref}
\newcommand{\ie}{i.e.\ }
\newcommand{\cf}{cf.\ }
\newcommand{\resp}{respectively\ }
\newcommand{\Subsection}{\subsection}
\providecommand{\nobreakdash}{\nobreak\hbox{-}}
\setlength{\emergencystretch}{2em}
\multlinegap0pt
\usepackage{bm}
\usepackage{bbm}
\usepackage{dsfont}
\usepackage{xspace}
\usepackage{cancel}
\usepackage{mathtools}
\usepackage{amsthm}
\makeatletter
\@ifundefined{mydef}{\newtheorem{mydef}{Mydef}}{}
\@ifundefined{myclaim}{\newtheorem{myclaim}{Claim}}{}
\@ifundefined{mylem}{\newtheorem{mylem}[mydef]{Lemma}}{}
\makeatother
\DeclareMathOperator{\osucc}{osucc}
\DeclareMathOperator{\stem}{stem}
\newcommand{\dL}{\mathbb{L}}
\newcommand{\uhr}{\upharpoonright}
\title[Total Failure of Approachability at Successors of Singulars ]{Total Failure of Approachability at Successors of Singulars of Countable Cofinality}
\author{Hannes Jakob}
\date{Source: arXiv:2508.02867v3 (CC BY 4.0)}
\begin{document}
\maketitle

\noindent\emph{Source and licence.} Total Failure of Approachability at Successors of Singulars of Countable Cofinality. Version arXiv:2508.02867v3 (CC BY 4.0), distributed under the Creative Commons Attribution 4.0 International licence, as recorded in the arXiv licence metadata for this exact version and shown on the abstract page https://arxiv.org/abs/2508.02867v3. Full rights notice: licenses/arxiv-2508.02867-CC-BY-4.0.txt.\par\smallskip

\noindent\textit{Editorial scope.} Counted unit: the Mylem whose source label is \texttt{LClosure} in the article (source0.tex lines 907--909, source label \texttt{LClosure}), delivered together with the article's own complete original proof of it (source0.tex lines 911--973, one \texttt{proof} environment of 63 lines). No mathematics is written, repaired or re-derived: statement and proof are the article's own text line for line. Required context retained uncounted: FubiniAssociative (lines 842--844); ProductTree (lines 874--876); ProductComplete (lines 895--897). Every definition, display and label of those lines is retained unchanged. Omitted from this package and not redistributed: 1--841, 845--873, 877--894, 898--906, 910--910, 974--1360. Nothing inside a retained excerpt is omitted or altered: every retained body is delivered exactly as the article's lines are written, so no omission receipt of any kind accompanies this package. Citations: the retained excerpts carry no citation command at all, so no bibliography entry of the article is transcribed and no reference is redistributed. Reference adaptation: none was needed and none was made; every reference inside the delivered unit resolves inside it, so no unresolved reference or citation is produced. Printed slips kept literal: no printed word, symbol, hypothesis or number of the counted statement or of its proof was altered. Layout and class: the article's own class is \texttt{amsart} with packages including babel, amssymb, mathrsfs, hyperref, xcolor, eucal, amsmath, amsthm, imakeidx, breqn, diagbox, verbatim, fancyhdr, tikz-cd; the wrapper is the lane's pinned \texttt{amsart} editorial wrapper at 10pt on A4, which restates the theorem-like environments the retained text needs and adds the article's own macro definitions that the retained text needs (\texttt{\textbackslash dL}, \texttt{\textbackslash osucc}, \texttt{\textbackslash stem}, \texttt{\textbackslash uhr}), with extra wrapper packages (bm, bbm, dsfont, xspace, cancel, mathtools). The delivered numbering is the unit's own. Assets: no retained span contains a figure environment, a graphics-inclusion command, a PDF-page inclusion, a table or a TikZ picture; no image file is copied into this package, the package declares no asset dependency and no image file is redistributed with it. Licence: version arXiv:2508.02867v3 of this article is distributed under the Creative Commons Attribution 4.0 International licence, as recorded in the arXiv licence metadata for this exact version and shown on the abstract page https://arxiv.org/abs/2508.02867v3. A full rights notice is delivered at licenses/arxiv-2508.02867-CC-BY-4.0.txt. Characters: every retained excerpt is pure ASCII, so the delivered engine, XeLaTeX with its default Latin Modern text font, reports no missing character.
	\begin{mylem}\label{FubiniAssociative}
		The Fubini product is associative. In other words, whenever $I,J,K$ are ideals (on $X,Y,Z$ respectively), then $(I\otimes J)\otimes K=I\otimes(J\otimes K)$.
	\end{mylem}

	\begin{mylem}\label{ProductTree}
		Let $J_0,\dots,J_{n-1}$ be ideals. Then a set $A\subseteq X_0\times\dots\times X_{n-1}$ is in $(J_0\otimes\dots\otimes J_{n-1})^+$ if and only if there is a $(J_k)_{k<n}$-Laver-tree $p$ such that $T(p)\subseteq A$.
	\end{mylem}

	\begin{mylem}\label{ProductComplete}
		Let $I$ and $J$ be ${<}\,\kappa$-complete. Then $I\otimes J$ is ${<}\,\kappa$-complete.
	\end{mylem}

	\begin{mylem}\label{LClosure}
		The poset $(\dL(\mu,(\kappa_n,I_n,B_n)_{n\in\omega}),\leq_0)$ is strongly ${<}\,\mu$-directed closed.
	\end{mylem}

	\begin{proof}
		Let $X\subseteq\dL(\mu,(\kappa_n,I_n,B_n)_{n\in\omega})$ be weakly directed with respect to $\leq_0$. Let $p:=\bigcap X$. Note that $p$ is not necessarily a condition in $\dL(\mu,(\kappa_n,I_n,B_n)_{n\in\omega})$, since $X$ is only weakly directed. E.g. it is possible that $X$ contains two conditions $q_0$ and $q_1$ such that for some $\alpha\in\osucc_{q_0}(\stem(q_0))\cap\osucc_{q_1}(\stem(q_1))$, the $\osucc$-sets of $\stem(q_i)^{\frown}\alpha$ are disjoint for $i=0,1$ (as long as this does not occur for too many $\alpha$). In that case, the intersection $q_0\cap q_1$ contains a terminal node, namely $\stem(q_0)^{\frown}\alpha$, and is therefore not a condition in $\dL$. However, we will show that $p$ can be refined to a condition in $\dL(\mu,(\kappa_n,I_n,B_n)_{n\in\omega})$ with the same stem. Let $s:=\stem(p)$ (which is also the stem of every $q\in X$). We will assume for notational simplicity that $s=\emptyset$.
		
		\setcounter{myclaim}{0}
		
		We first prove a general claim. Let $\dL:=\dL(\mu,(\kappa_n,I_n,B_n)_{n\in\omega})$.
		
		\begin{myclaim}\label{LClosureClaim1}
			If $q_0,q_1\in\dL$, $\stem(q_0)=\stem(q_1)=\emptyset$ and $q_0\Vdash\check{q}_1\in\Gamma_{\dL}$, then for any $n\in\omega$,
			$$\osucc_{q_0}^n(\emptyset)\subseteq^*\osucc_{q_1}^n(\emptyset)$$
			where $\subseteq^*$ is defined with regard to $I_0\otimes\dots\otimes I_{n-1}$.
		\end{myclaim}
		
		\begin{proof}
			Assume toward a contradiction that $A:=\osucc_{q_0}^n(\emptyset)\smallsetminus\osucc_{q_1}^n(\emptyset)$ is $ (I_0\otimes\dots\otimes I_{n-1})$-positive. By Lemma \ref{ProductTree} there is an $(I_k)_{k<n}$-tree $T$ such that $[T]\subseteq A$. Let $q':=\{s\in q_0\;|\;s\uhr n\in T\}$. We claim that $q'$ is a condition in $\dL$: Let $s\in q'$. If $|s|\geq n$, then $\osucc_{q'}(s)=\osucc_{q_0}(s)\in I_{|s|}^+$ so we are done. If $|s|<n$, then $\osucc_{q'}(s)=\osucc_T(s)$, since $[T]\subseteq A\subseteq\osucc_{q_0}^n(\emptyset)$. But that set is in $I_{|s|}^+$ as well by assumption. It follows that $q'\leq_0q_0$ and $q'\perp q_1$, since for every $s\in q'$ with $|s|=n$, $s\in A$ and therefore $s\notin\osucc_{q_1}^n(\emptyset)$. This  contradicts the fact that $q_0\Vdash\check{q}_1\in\Gamma_{\dL}$.
		\end{proof}
		
		This claim implies that the first $\osucc$-set is as required:
		
		\begin{myclaim}\label{LClosureClaim2}
			$\osucc_p(\emptyset)\in I_0^+$.
		\end{myclaim}
		
		\begin{proof}
			We claim that the collection $\{\osucc_q(\emptyset)\;|\;q\in X\}$ is a weakly directed subset of $(B_0,\subseteq)$. Given $q_0,q_1\in X$, let $q_2\in X$ be such that $q_2\Vdash q_0,q_1\in\Gamma_{\dL}$. It follows from Claim \ref{LClosureClaim1} with $n=0$ that $\osucc_{q_2}(\emptyset)\subseteq^*\osucc_{q_0}(\emptyset)\cap\osucc_{q_1}(\emptyset)$, so $\osucc_{q_2}(\emptyset)$ witnesses that $\osucc_{q_0}(\emptyset)$ and $\osucc_{q_1}(\emptyset)$ are weakly compatible.
			
			Since $(B_0,\subseteq)$ is strongly ${<}\,\mu$-directed closed and since $\osucc_p(\emptyset)=\bigcap_{q\in X}\osucc_q(\emptyset)$, the latter set is in $I_0^+$.
		\end{proof}
		
		We also have that at least for ``most'' $t\in\osucc_p^n(\emptyset)$, the $\osucc$-sets are large enough:
		
		\begin{myclaim}\label{LClosureClaim3}
			For $n\in\omega$, $\{t\in\osucc_p^n(\emptyset)\;|\;\osucc_p(t)\in I_n\}\in I_0\otimes\dots\otimes I_{n-1}$.
		\end{myclaim}
		
		\begin{proof}
			Let $q_0,q_1\in X$. Then there is $q_2\in X$ such that $q_2\Vdash q_0,q_1\in\Gamma_{\dL}$. By Claim \ref{LClosureClaim1} we have that $\osucc_{q_2}^{k}(\emptyset)\subseteq^*\osucc_{q_0}^{k}(\emptyset),\osucc_{q_1}^{k}(\emptyset)$. Let $i<2$ be arbitrary. Then $\osucc_{q_2}^{n+1}\smallsetminus\osucc_{q_i}^{n+1}\in I_0\otimes\dots\otimes I_{n}$ and therefore
			$$\{t\in\osucc_{q_2}^n(\emptyset)\;|\;\{x\in[\kappa_{n-1}]^{<\mu^+}\;|\;(t,x)\in\osucc_{q_2}^{n+1}(\emptyset)\smallsetminus\osucc_{q_i}^{n+1}(\emptyset)\}\in I_n^+\}$$
			is in $I_0\otimes\dots\otimes I_{n-1}$
			Ergo, we have
			$$\{t\in\osucc_{q_2}^n(\emptyset)\;|\;\osucc_{q_2}(t)\not\subseteq^*\osucc_{q_i}(t)\}\in I_0\otimes\dots\otimes I_{n-1}.$$
			It follows that
			$$A_{q_0,q_1,q_2}:=\{t\in\osucc_{q_2}^n(s)\;|\;\exists i<2(t\notin\osucc_{q_i}(\emptyset)\vee\osucc_{q_2}(t)\not\subseteq^*\osucc_{q_i}(t))\}$$
			is in $I_0\otimes\dots\otimes I_{n-1}$ since it is the union of four sets in $I_0\otimes\dots\otimes I_{n-1}$. Let $A$ be the union of all $A_{q_0,q_1,q_2}$ where $q_0,q_1\in X$ and $q_2$ witnesses their weak compatibility. Then $A\in I_0\otimes\dots\otimes I_{n-1}$ since that ideal is ${<}\,\mu$-complete by Lemma \ref{ProductComplete}. It follows that whenever $t\in\osucc_p^n(\emptyset)\smallsetminus A$, the collection $\{\osucc_q(t)\;|\;q\in X\}$ is weakly directed (first note that $t\in q$ whenever $q\in X$): Given $q_0,q_1\in X$, let $q_2$ witness their weak compatibility. By assumption $t\notin A_{q_0,q_1,q_2}$ and so $\osucc_{q_2}(t)\subseteq^*\osucc_{q_0}(t),\osucc_{q_1}(t)$. So whenever $t\in\osucc_p^n(\emptyset)\smallsetminus A$, we have that $\osucc_p(t)=\bigcap_{q\in X}\osucc_q(t)\in I_{|t|}^+$, since $(B_n,\subseteq)$ is strongly ${<}\,\mu$-directed closed. It follows that
			$$\{t\in\osucc_p^n(\emptyset)\;|\;\osucc_p(t)\in I_n\}\subseteq A\in I_0\otimes\dots\otimes I_{n-1}.$$
		\end{proof}
		
		Now we let $q'$ consist of all $t\in p$ such that $\osucc_p(t)\in I_{|t|}^+$ and for every $n$,
		$$\{r\in\osucc_p^n(t)\;|\;\osucc_p^n(t^{\frown}r)\in I_{|t|+n}\}\in I_{|t|}\otimes\dots\otimes I_{|t|+(n-1)}.$$
		Let $q$ consist of all those $t\in q'$ such that every initial segment of $t$ is also in $q'$. We claim that $q\in\dL$ and $\stem(q)=\emptyset$. By Claim \ref{LClosureClaim2} and Claim \ref{LClosureClaim3}, $\emptyset\in q$ and so $q$ is nonempty. It is clearly closed under initial segments by construction. So we have to show that it splits widely enough. Let $t\in q$. By assumption, $\osucc_p(t)\in I_{|t|}^+$. Additionally, $\osucc_q(t)$ consists of all $x\in\osucc_p(t)$ such that $\osucc_p(t^{\frown}x)\in I_{|t|+1}^+$ and for every $n$,
		$$\{r\in\osucc_p^n(t^{\frown}x)\;|\;\osucc_p(t^{\frown}x^{\frown}r)\in I_{|t|+1+n}\}\in I_{|t|+1}\otimes\dots\otimes I_{|t|+n}.$$
		So suppose $\osucc_q(t)\in I_{|t|}$. Then $\osucc_p(t)\smallsetminus\osucc_q(t)\in I_{|t|}^+$ and so for every $x\in\osucc_p(t)\smallsetminus\osucc_q(t)$, either $\osucc_p(t^{\frown}x)\in I_{|t|+1}$ or there is some $n$ such that
		$$\{r\in\osucc_p^n(t^{\frown}x)\;|\;\osucc_p(t^{\frown}x^{\frown}r)\in I_{|t|+1+n}\}\in (I_{|t|+1}\otimes\dots\otimes I_{|t|+n})^+.$$
		By the ${<}\,\mu^+$-completeness of $I_{|t|}$, one of those cases has to occur for $I_{|t|}^+$ many $x$. If it is the first one, we directly have
		$$\{x\in\osucc_p(t)\;|\;\osucc_p(t^{\frown}x)\in I_{|t|+1}\}\in I_{|t|}^+$$
		contradicting the fact that $t\in q$, or more precisely, the second requirement for $n=1$. In the second case, we have that the set
		$$\{r\in\osucc_p^n(t^{\frown}x)\;|\;\osucc_p(t^{\frown}x^{\frown}r)\in I_{|t|+1+n}\}$$
		is in $(I_{|t|+1}\otimes\dots\otimes I_{|t|+n})^+$. Let $A:=\{r\in\osucc_p^{n+1}(t)\;|\;\osucc_p(t^{\frown}r)\in I_{|t|+1+n}\}$. Then the above equation states that
		$$\{x\in [\kappa_{|t|}]^{<\mu^+}\;|\;\{r\in\prod_{i\in[|t|+1,|t|+n]}[\kappa_i]^{<\mu^+}\;|\;(x,r)\in A\}\in(I_{|t|}\otimes\dots\otimes I_{|t|+n})^+\}\in I_{|t|}^+.$$
		In other words, $A\in(I_{|t|}\otimes(I_{|t|+1}\otimes\dots\otimes I_{|t|+n}))^+=(I_{|t|}\otimes\dots\otimes I_{|t|+n})^+$, where the equality follows from Lemma \ref{FubiniAssociative}. But this clearly contradicts the fact that $t\in q$, or more precisely, the second requirement for $n+1$.
		
		In summary, $q$ is a condition in $\dL$ and a $\leq_0$-lower bound of $X$. This finishes the proof.
	\end{proof}

\end{document}
